\documentclass[12pt]{article}
\usepackage[utf8]{inputenc}
\usepackage[T1]{fontenc}
\usepackage[margin=1in]{geometry}
\usepackage{setspace}
\usepackage{xcolor}
\usepackage[colorlinks=true, urlcolor=blue!60!black, linkcolor=blue!60!black]{hyperref}
\setlength{\parindent}{2em}
\setlength{\parskip}{1.25em}
\renewcommand{\baselinestretch}{1.0}
\begin{document}
\begin{center}
{\LARGE \textit{Seminole Tribe of Florida v.\ Florida}, 517 U.S.\ 44 (1996)\par}
\vspace{0.8em}
{\large Case Brief --- Four-Page Edition\par}
\end{center}
\vspace{1em}

\section{Memory Jogger}

Overruling \textit{Pennsylvania v.\ Union Gas}, the Court held 5--4 that Congress cannot use its Article I powers, including the Indian Commerce Clause, to abrogate state sovereign immunity, and that \textit{Ex parte Young} was unavailable to enforce the Indian Gaming Regulatory Act's detailed remedial scheme against Florida's Governor (\href{https://www.law.cornell.edu/supct/html/94-12.ZO.html}{Cornell LII}).

\section{Detailed Case Facts}

Under the Indian Commerce Clause, Congress enacted the Indian Gaming Regulatory Act (IGRA) in 1988. Class III gaming (slot machines, casino games, and the like) is lawful only under a Tribal-State compact, and IGRA required a State, upon a tribe's request, to ``negotiate with the Indian tribe in good faith to enter into such a compact.'' To enforce that duty, IGRA let tribes sue States in federal court; if the court found bad faith, it could order a compact within 60 days, then a mediator would choose between each side's last best offer,'' and if the State still refused, the Secretary of the Interior would prescribe gaming procedures. In 1991 the Seminole Tribe sued Florida and Governor Lawton Chiles, alleging that they refused to negotiate.

\section{Procedural History}

The district court denied Florida's motion to dismiss on Eleventh Amendment grounds, holding that Congress validly abrogated immunity. On interlocutory appeal, the Eleventh Circuit reversed, holding that the Indian Commerce Clause did not authorize abrogation and that \textit{Ex parte Young} did not permit suit against the Governor. The Supreme Court granted certiorari and affirmed on March 27, 1996.

\section{Judicial Votes}

Affirmed, \textbf{5--4}.
\begin{itemize}
\item \textbf{Majority:} Chief Justice Rehnquist, joined by Justices O'Connor, Scalia, Kennedy, and Thomas.
\item \textbf{Dissenting:} Justice Stevens; Justice Souter, joined by Justices Ginsburg and Breyer.
\end{itemize}

\section{Holding}

To abrogate state immunity, Congress must unequivocally express its intent and act under a valid grant of constitutional authority. IGRA satisfied the first requirement, but \textit{Union Gas} is overruled: the Eleventh Amendment restricts Article III judicial power, and Article I powers cannot abrogate the immunity of unconsenting States. Congress may still abrogate under Section~5 of the Fourteenth Amendment (\textit{Fitzpatrick v.\ Bitzer}), which postdates and altered the Eleventh. Where Congress has prescribed a detailed remedial scheme, courts may not permit an \textit{Ex parte Young} suit that circumvents it. The breadth of the Article I bar was later qualified by \textit{Katz}, \textit{PennEast}, and \textit{Torres}.

\section{Analysis of Opinions}

\textbf{Majority (Rehnquist).} Since \textit{Hans v.\ Louisiana}, the Eleventh Amendment has confirmed a broader ``presupposition'' that States are not amenable to private suit without consent. IGRA's references to ``the State'' as defendant were ``unmistakably clear.'' Seeking injunctive rather than monetary relief did not matter, because immunity guards against ``the indignity of subjecting a State to the coercive process of judicial tribunals at the instance of private parties.'' \textit{Fitzpatrick} rested on the Fourteenth Amendment's later adoption, while \textit{Union Gas} was a fractured plurality with no principled distinction between the Interstate and Indian Commerce Clauses. Rehnquist dismissed Souter's textualism as a ``straw man.'' On \textit{Ex parte Young}, analogizing to \textit{Schweiker v.\ Chilicky}, an officer suit would expose the Governor to full contempt powers and render IGRA's modest scheme superfluous.

\textbf{Stevens (dissenting).} The ruling was a ``sharp break with the past'' and an ``affront to a coequal branch,'' leaving victims of state violations of copyright, patent, bankruptcy, and environmental laws without remedy. \textit{Chisholm}'s Iredell dissent and \textit{Hans} treated immunity as a common-law presumption Congress could displace.

\textbf{Souter, with Ginsburg and Breyer (dissenting).} The Eleventh Amendment merely repealed the Citizen-State Diversity Clauses, leaving federal-question jurisdiction untouched. \textit{Hans} never made immunity immune from abrogation, and elevating an unwritten common-law rule above Article I repeated the error of \textit{Lochner}. \textit{Young} is ``indispensable to the establishment of constitutional government and the rule of law,'' and IGRA's procedures could govern a \textit{Young} suit.

\section{Examples: Future Applications}

\begin{enumerate}
\item \textbf{Same-Side: Patent and Copyright Damages.} Congress authorizes infringement suits against States under the Intellectual Property Clause. Article I cannot abrogate, as \textit{Florida Prepaid} and \textit{Allen v.\ Cooper} confirmed.
\item \textbf{Same-Side: Age Discrimination.} Under the Commerce Clause and Section~5, Congress lets state employees sue for age discrimination. Without a record of unconstitutional discrimination, abrogation fails (\textit{Kimel}).
\item \textbf{Opposite-Side: Family Leave.} Congress authorizes damages for family-leave violations, backed by evidence of gender stereotyping. Section~5 supports abrogation (\textit{Nevada v.\ Hibbs}).
\item \textbf{Opposite-Side: \textit{Ex parte Young} Suit.} A carrier seeks to enjoin state commissioners from enforcing an order violating a federal statute with only a general review provision. \textit{Young} remains available (\textit{Verizon Maryland}).
\item \textbf{Fence-Sitter: Bankruptcy or War Powers.} \textit{Seminole Tribe} said Article I cannot circumvent federal-jurisdiction limits but acknowledged ``surrender of this immunity in the plan of the convention''; later cases found such surrender for bankruptcy, eminent domain, and military reemployment.
\end{enumerate}

\section{Critique}

\textbf{Originalist views.} Critics note the irony of Justices Scalia and Thomas joining an opinion disparaging ``blind reliance upon the text'' for an atextual presupposition; following Souter, Akhil Amar, and William Fletcher, they read the Amendment to bar only diversity suits. William Baude and Stephen Sachs defend the result: immunity was a preexisting general-law backdrop that neither Article III nor Article I authorized Congress to abrogate.

\textbf{Progressive views.} Vicki Jackson and Erwin Chemerinsky argue that the case forced Congress to rely on Section~5 just as \textit{City of Boerne} constricted it, creating a federalism pincer, and that curtailing \textit{Ex parte Young} threatens state compliance with federal law. Federal agency enforcement is an unrealistic substitute for private suits.

\section{Key Quotations}

\begin{quote}
``Even when the Constitution vests in Congress complete lawmaking authority over a particular area, the Eleventh Amendment prevents congressional authorization of suits by private parties against unconsenting States. The Eleventh Amendment restricts the judicial power under Article III, and Article I cannot be used to circumvent the constitutional limitations placed upon federal jurisdiction.''
\end{quote}
\noindent --- Chief Justice Rehnquist, Opinion of the Court.

\begin{quote}
``[T]herefore, where Congress has prescribed a detailed remedial scheme for the enforcement against a State of a statutorily created right, a court should hesitate before casting aside those limitations and permitting an action against a state officer based upon \textit{Ex parte Young}.''
\end{quote}
\noindent --- Chief Justice Rehnquist, Opinion of the Court.

\begin{quote}
``The majority's opinion does not simply preclude Congress from establishing the rather curious statutory scheme under which Indian tribes may seek the aid of a federal court to secure a State's good faith negotiations over gaming regulations. Rather, it prevents Congress from providing a federal forum for a broad range of actions against States, from those sounding in copyright and patent law, to those concerning bankruptcy, environmental law, and the regulation of our vast national economy.''
\end{quote}
\noindent --- Justice Stevens, dissenting.

\end{document}
